% =====================================================================
% Capítulo 2 (continuación): Secs. 2.7-2.8
%   Pruebas de hipótesis para tablas de conteo y SigTree
% =====================================================================

% ---------------------------------------------------------------------
\section{Pruebas de hipótesis para tablas de conteo}\label{sec:pruebas}

La Sección~\ref{sec:idiv-poisson} mostró que la interacción que queda dentro de un
bloque se mide con el estadístico $G^2$ de una prueba de independencia. Para usar
ese estadístico como criterio de refinamiento hay que resolver tres cuestiones. La
primera es cómo obtener una distribución de referencia válida cuando las tablas
son pequeñas o escasas y cuando los bloques se eligieron con los mismos datos. La
segunda es cómo controlar el error al realizar muchas pruebas, una por cada bloque
de un árbol de refinamientos. La tercera es cómo evitar que, con millones de
viajes, cualquier diferencia mínima resulte significativa. Esta sección presenta
las herramientas estadísticas que responden a cada una.

A lo largo de la sección se considera una tabla de contingencia genérica
$x = (x_{ij})$ de $a$ filas y $b$ columnas, con totales $x_{i+}$, $x_{+j}$ y
$x_{++}$. En la propuesta, esa tabla será la restricción de la matriz de
co-ocurrencia a un bloque.

\subsection{Pruebas de independencia}\label{sec:est-independencia}

La hipótesis nula de independencia afirma que la probabilidad de cada celda es el
producto de las probabilidades marginales, $\pi_{ij} = \pi_{i+}\,\pi_{+j}$, o, en
el modelo de Poisson, que las medias se factorizan como $\mu_{ij} = \alpha_i
\beta_j$. En cualquiera de los dos casos, los conteos esperados estimados bajo la
hipótesis nula son
\begin{equation}\label{eq:esperados}
  \hat{E}_{ij} = \frac{x_{i+}\, x_{+j}}{x_{++}} .
\end{equation}
Las dos pruebas clásicas comparan los conteos observados con estos esperados
\cite{Agresti2013}. La prueba de Pearson usa el estadístico
\begin{equation}\label{eq:pearson}
  X^2 = \sum_{i,j} \frac{(x_{ij} - \hat{E}_{ij})^2}{\hat{E}_{ij}} ,
\end{equation}
y la prueba de razón de verosimilitud usa
\begin{equation}\label{eq:g2}
  G^2 = 2 \sum_{i,j} x_{ij} \log \frac{x_{ij}}{\hat{E}_{ij}} .
\end{equation}
Bajo la hipótesis nula, y cuando los conteos esperados son grandes, ambos
estadísticos siguen aproximadamente una distribución $\chi^2$ con
$(a-1)(b-1)$ grados de libertad, y son asintóticamente equivalentes
\cite{Agresti2013}. Con conteos esperados pequeños la aproximación se deteriora;
la regla práctica más difundida pide que la mayoría de los conteos esperados sean
al menos cinco.

El estadístico $G^2$ tiene dos propiedades que lo hacen preferible en esta tesis.
La primera es su conexión con el marco de Bregman: como los esperados
\eqref{eq:esperados} reproducen el total de la tabla, los términos lineales de la
I-divergencia se cancelan y $G^2 = 2\,D_I(x \,\|\, \hat{E})$. Por la
ecuación~\eqref{eq:g2-im}, esto equivale a $G^2 = 2\,x_{++}\,\hat{I}(R;T)$: el
estadístico es la información mutua de la tabla escalada por su masa. La segunda
es que $G^2$ admite una partición exacta: cuando las filas de una tabla se
agrupan, el $G^2$ de la tabla completa es igual al $G^2$ de la tabla colapsada más
la suma de los $G^2$ de las subtablas de cada grupo \cite{Agresti2013}. El
estadístico de Pearson no cumple esta identidad de forma exacta. Esta propiedad
es la que permite, en el Capítulo~\ref{chap:propuesta}, repartir la interacción
de un bloque entre la parte que explica un corte y la que queda dentro de los
hijos.

\subsection{Pruebas condicionales y p-valores por simulación}\label{sec:est-permutacion}

La distribución $\chi^2$ es solo una aproximación. Para tablas con celdas
escasas, que son frecuentes en bloques de poca actividad, conviene una prueba cuya
validez no dependa del tamaño de la muestra. La idea es eliminar los parámetros
desconocidos condicionando sobre sus estadísticos suficientes. Bajo la hipótesis
nula de independencia, tanto en el modelo de Poisson como en el multinomial, la
distribución de la tabla condicionada a sus totales de fila y de columna es la
hipergeométrica multivariada, que no depende de los parámetros $\alpha_i$ y
$\beta_j$ \cite{Agresti2013}. El p-valor exacto es la probabilidad, bajo esa
distribución, de obtener un estadístico al menos tan extremo como el observado.
Para tablas $2 \times 2$ esta es la prueba exacta de Fisher.

Enumerar todas las tablas con los mismos totales es inviable salvo en tablas muy
pequeñas, por lo que el p-valor se estima por simulación. El algoritmo de
\cite{Patefield1981} genera tablas aleatorias con totales de fila y columna
fijos, con probabilidad igual a la hipergeométrica multivariada. Con $B$ tablas
simuladas y sus estadísticos $t_1, \dots, t_B$, el p-valor se estima como
\begin{equation}\label{eq:pval-mc}
  \hat{p} = \frac{1 + \#\{k : t_k \ge t_{\mathrm{obs}}\}}{1 + B} .
\end{equation}
\cite{Phipson2010} muestran que sumar uno en el numerador y en el denominador no
es un ajuste arbitrario: la tabla observada es, bajo la hipótesis nula, una
realización más de la misma distribución, y la fórmula~\eqref{eq:pval-mc} produce
una prueba exacta, mientras que la proporción sin corregir $\#\{k : t_k \ge
t_{\mathrm{obs}}\}/B$ puede valer cero y subestima el p-valor. Una consecuencia
práctica es que el menor p-valor alcanzable es $1/(B+1)$, así que para probar a un
nivel $\alpha$ se necesita $B \ge 1/\alpha - 1$ como mínimo.

Cuando se realizan muchas pruebas, la mayoría con p-valores grandes, simular $B$
tablas en todas ellas desperdicia cómputo. \cite{Besag1991} propusieron un
procedimiento secuencial: se simulan tablas hasta que $h$ de ellas superan al
estadístico observado, o hasta alcanzar un máximo $B$. Si la simulación se detiene
tras $l$ tablas, el p-valor es $h/l$; si no, se usa la fórmula anterior. El
procedimiento mantiene la validez de la prueba y detiene pronto la simulación en
los casos claramente no significativos.

Un aspecto de las pruebas por simulación que se usará después es que el
estadístico puede ser cualquier función de la tabla, incluida una que involucre
una búsqueda. Si $t$ es, por ejemplo, el máximo de $G^2$ sobre un conjunto de
divisiones candidatas, el p-valor~\eqref{eq:pval-mc} sigue siendo válido siempre
que en cada tabla simulada se repita la misma búsqueda. Esta es la base del ajuste
por selección de la Sección~\ref{sec:est-seleccion}.

\subsection{Comparaciones múltiples}\label{sec:est-multiplicidad}

Cuando se realizan $M$ pruebas, cada una a un nivel $\alpha$, la probabilidad de
cometer al menos un falso rechazo crece con $M$. Dos medidas de error resumen el
problema. La tasa de error por familia (FWER) es la probabilidad de cometer al
menos un falso rechazo. La tasa de falsos descubrimientos (FDR) es la proporción
esperada de falsos rechazos entre todos los rechazos \cite{Benjamini1995}.
Controlar la FWER es más estricto; controlar la FDR tolera algunos falsos
rechazos a cambio de más potencia, lo que suele ser preferible en análisis
exploratorios con muchas hipótesis.

La corrección de Bonferroni prueba cada hipótesis al nivel $\alpha/M$ y controla la
FWER sin ningún supuesto sobre la dependencia entre pruebas. Su costo es la
potencia: cuando las pruebas están correlacionadas, como ocurre al evaluar
divisiones candidatas de un mismo bloque, que comparten casi todos sus datos, la
corrección es mucho más conservadora de lo necesario. El ajuste \emph{min-P} de
\cite{Westfall1993} resuelve este problema usando remuestreo: en lugar de dividir
$\alpha$ entre $M$, compara el menor p-valor observado, o de forma equivalente el
mayor estadístico, con su distribución bajo la hipótesis nula conjunta, obtenida
por simulación. Así el ajuste tiene en cuenta a la vez el número de pruebas y la
correlación entre ellas. \cite{Benjamini1995} propusieron, para la FDR, un
procedimiento secuencial sobre los p-valores ordenados que controla esa tasa
cuando las pruebas son independientes.

\paragraph{Pruebas jerárquicas.} En muchos problemas las hipótesis forman un
árbol, y la hipótesis de un nodo solo tiene sentido si la de su padre fue
rechazada. En ese caso no es necesario corregir por todas las pruebas posibles,
sino solo por las que efectivamente se realizan. \cite{Meinshausen2008} propuso un
procedimiento de arriba abajo para grupos jerárquicos de variables, en el que el
nivel asignado a cada grupo es proporcional a su tamaño, y mostró que controla la
FWER sobre toda la jerarquía. \cite{Kimes2017} aplicaron una idea análoga al
agrupamiento jerárquico: recorren el dendrograma desde la raíz, prueban en cada
nodo si la división en dos grupos es significativa mediante una prueba por
simulación, y controlan la FWER asignando a cada nodo un umbral proporcional a su
número de observaciones. Para la FDR, \cite{Yekutieli2008} propuso aplicar el
procedimiento de \cite{Benjamini1995} dentro de cada familia de hijos de una
hipótesis rechazada, y acotó la FDR resultante sobre el árbol completo.

Estos procedimientos aprovechan una estructura que también tiene el problema de
esta tesis. Si la hipótesis nula de independencia se cumple en un bloque, se
cumple en todos sus sub-bloques, porque cualquier subtabla de una tabla con medias
factorizadas también las tiene factorizadas. Por lo tanto, las hipótesis nulas
verdaderas forman subárboles completos, y un falso rechazo solo puede producirse
por primera vez en la raíz de uno de ellos.

\subsection{Inferencia después de la selección}\label{sec:est-seleccion}

Las garantías anteriores suponen que las hipótesis se fijan antes de ver los
datos. En el agrupamiento no ocurre así: los grupos, y por tanto las hipótesis que
se prueban sobre ellos, se eligen con los mismos datos. Este problema se conoce
como \emph{doble uso de los datos}. \cite{Gao2024} mostraron que, cuando los grupos
se definen mediante agrupamiento, una prueba clásica de diferencia de medias
tiene una tasa de error de tipo~I extremadamente inflada, y que el problema
persiste aunque se usen dos conjuntos de datos independientes para definir los
grupos y para probarlos. La razón es que el procedimiento de agrupamiento
encuentra grupos distintos incluso cuando no hay estructura, y la prueba clásica
no lo tiene en cuenta.

Existen tres formas de abordar el problema, que la Sección~\ref{sec:ea_inferencia}
revisa con más detalle:

\begin{enumerate}
  \item \textbf{Incluir la selección en la simulación.} Si la prueba es por
  simulación, basta con repetir el procedimiento de selección completo en cada
  réplica nula, como se indicó al final de la Sección~\ref{sec:est-permutacion}.
  El p-valor resultante es válido para la hipótesis nula global de que no hay
  estructura en la tabla, aunque no para hipótesis más específicas.
  \item \textbf{Inferencia selectiva.} Se calcula el p-valor condicionando sobre
  el evento de que el procedimiento haya producido los grupos observados. Es el
  enfoque de \cite{Gao2024} para el agrupamiento jerárquico y de
  \cite{Watanabe2021si} para modelos de bloques latentes. Ofrece garantías
  formales, pero se ha desarrollado sobre todo para datos gaussianos y su costo
  computacional es alto.
  \item \textbf{División de cada observación.} Para conteos de Poisson,
  \cite{Neufeld2024cs} proponen el \emph{count splitting}: cada conteo se divide
  como $x_{ij} = x^{(1)}_{ij} + x^{(2)}_{ij}$, con $x^{(1)}_{ij} \mid x_{ij} \sim
  \mathrm{Binomial}(x_{ij}, \varepsilon)$. Si $x_{ij} \sim
  \mathrm{Poisson}(\mu_{ij})$, las dos partes son independientes y de Poisson, con
  medias $\varepsilon\mu_{ij}$ y $(1-\varepsilon)\mu_{ij}$. La estructura se
  elige con $x^{(1)}$ y se prueba con $x^{(2)}$ como si hubiera sido fijada de
  antemano. \cite{Neufeld2024dt} generalizan la idea, con el nombre de
  \emph{data thinning}, a otras familias de distribuciones cerradas bajo
  convolución. Si los conteos presentan sobredispersión, las dos partes dejan de
  ser independientes y la garantía se debilita.
\end{enumerate}

\subsection{Tamaño de efecto y sobredispersión}\label{sec:est-efecto}

\paragraph{Tamaño de efecto.} Con muestras grandes, una prueba de hipótesis
detecta desviaciones de la hipótesis nula por pequeñas que sean. En el caso de la
independencia esto se ve directamente: por la relación $G^2 = 2\,x_{++}\,\hat
I(R;T)$, si la información mutua de la población es positiva, aunque sea
mínima, el estadístico crece de forma proporcional al número de observaciones y
el p-valor tiende a cero. Los propios \cite[Sec.~5]{Hu2025} identifican este
fenómeno como una limitación de SigTree. Con millones de viajes, la
significancia deja de distinguir entre interacciones relevantes e irrelevantes.

La respuesta habitual es acompañar la significancia con una medida de tamaño de
efecto que no crezca con la muestra. Para tablas de contingencia hay dos opciones
naturales. La primera es la propia información mutua por observación,
$\hat I(R;T) = G^2/(2\,x_{++})$, medida en nats, que conecta directamente con la
función objetivo del co-agrupamiento. La segunda es la $V$ de Cramér,
\begin{equation}\label{eq:cramer}
  V = \sqrt{\frac{X^2}{x_{++}\,\min(a-1,\, b-1)}} ,
\end{equation}
que toma valores entre 0, bajo independencia, y 1, bajo asociación perfecta
\cite{Agresti2013}. Ambas permiten fijar un umbral de relevancia independiente
del tamaño de la muestra.

\paragraph{Sobredispersión.} El modelo de Poisson impone que la varianza de cada
conteo sea igual a su media. En datos reales es frecuente que la varianza sea
mayor, $\mathrm{Var}(x_{ij}) = \varphi\,\mu_{ij}$ con $\varphi > 1$, porque los
eventos no son independientes: en movilidad, la actividad de una celda varía de un
día a otro por el día de la semana, el clima o eventos puntuales. En ese caso los
estadísticos $X^2$ y $G^2$ se inflan aproximadamente por el factor $\varphi$, y
una prueba calibrada bajo el modelo de Poisson rechaza la hipótesis nula con
demasiada frecuencia. El enfoque de cuasi-verosimilitud corrige este efecto
dividiendo los estadísticos entre una estimación de $\varphi$, por ejemplo el
cociente entre el estadístico de Pearson y sus grados de libertad
\cite{McCullagh1989}. Cuando la matriz se construye agregando varios días, como en
esta tesis, $\varphi$ también puede estimarse a partir de la variación de cada
celda entre días.

% ---------------------------------------------------------------------
\section{Agrupamiento guiado por pruebas de hipótesis: SigTree}\label{sec:sigtree}

Los árboles de decisión son uno de los modelos más usados para producir
agrupamientos interpretables, porque cada grupo queda descrito por la secuencia de
reglas que lleva a su hoja \cite{Hu2025}. En la mayoría de los métodos, sin
embargo, la decisión de dividir un nodo no tiene una interpretación estadística.
La idea de usar pruebas de hipótesis para decidir las divisiones tiene un
antecedente supervisado en CHAID \cite{Kass1980}, que elige las divisiones de un
árbol de clasificación con pruebas $\chi^2$ entre cada división candidata y la
variable de clase. \cite{Hu2025} trasladaron esta idea al caso no supervisado con
SigTree, un algoritmo que construye un árbol de decisión binario para agrupar
datos categóricos y que usa una prueba de hipótesis tanto para elegir cada
división como para decidir cuándo detenerse.

\subsection{Divisiones candidatas y pseudo-clases}

Sea un conjunto de $n$ muestras descritas por $m$ atributos categóricos, donde el
atributo $q$ tiene $k_q$ categorías. En cada nodo, SigTree considera como división
candidata cada categoría $c$ de cada atributo $j$: la división separa las muestras
que toman el valor $c$ en el atributo $j$ de las que no, y se representa con una
variable binaria $Z = \mathbb{1}[x_j = c]$. Como no hay una variable de clase,
cada uno de los otros $m-1$ atributos actúa como \emph{pseudo-clase}. La hipótesis
nula de cada prueba es que $Z$ es independiente del atributo $q$, y la hipótesis
nula global es que $Z$ es independiente de todos los demás atributos. La intuición
es que, si las muestras del nodo no tienen estructura de grupos, ninguna división
candidata estará asociada con la mayoría de las pseudo-clases \cite[Sec.~3]{Hu2025}.

\subsection{Estadístico y criterio de división}

Para cada par formado por la división $Z$ y un atributo $q \ne j$ se construye la
tabla de contingencia $2 \times k_q$ de sus frecuencias conjuntas y se calcula el
estadístico de Pearson $X^2(Z, q)$ de la ecuación~\eqref{eq:pearson}. En tablas
$2 \times 2$ con algún conteo esperado menor que cinco se aplica la corrección de
Yates. Los estadísticos de todos los atributos se suman,
\begin{equation}\label{eq:sigtree}
  X^2(Z) = \sum_{q \ne j} X^2(Z, q), \qquad
  \mathrm{gl}(Z) = \sum_{q \ne j} (k_q - 1),
\end{equation}
y el p-valor de la división se obtiene de la distribución $\chi^2$ con
$\mathrm{gl}(Z)$ grados de libertad, bajo el supuesto de que las pruebas
individuales son independientes \cite[Sec.~3.3]{Hu2025}. En cada nodo se elige la
división candidata con menor p-valor.

\subsection{Criterios de parada}

La división elegida solo se aplica si su p-valor es menor o igual que un nivel
ajustado $\alpha^\ast$. Para controlar la FWER al elegir la mejor de muchas
candidatas, SigTree usa la corrección de Bonferroni con un divisor constante,
el número total de categorías de todos los atributos, $|Q|$:
\begin{equation}\label{eq:sigtree-alpha}
  \alpha^\ast = \frac{\alpha}{|Q|} ,
\end{equation}
que se aplica por igual en todos los nodos \cite[Sec.~3.4]{Hu2025}. Además, se
usan dos reglas triviales: no se dividen nodos con cinco muestras o menos, y no se
aceptan divisiones que produzcan un hijo con menos de tres muestras. El
crecimiento termina cuando ninguna hoja admite una división significativa, y las
hojas del árbol son los grupos finales. El resultado es un algoritmo con un solo
parámetro con interpretación estadística clara, el nivel $\alpha$, que no
requiere fijar el número de grupos de antemano.

\subsection{Limitaciones}

Los autores reconocen tres limitaciones \cite[Sec.~5]{Hu2025}. Evaluar todas las
divisiones candidatas en cada nodo es costoso. El algoritmo no optimiza una
función objetivo explícita, por lo que no hay garantía de que el árbol obtenido
sea el mejor posible. Y, como se discutió en la Sección~\ref{sec:est-efecto}, el
estadístico $\chi^2$ crece con el número de muestras y de atributos, lo que empuja
los p-valores hacia cero y dificulta controlar los falsos descubrimientos con el
umbral~\eqref{eq:sigtree-alpha}. A estas se pueden añadir dos observaciones desde
la Sección~\ref{sec:pruebas}: el p-valor de la división elegida se calcula como si
esa división se hubiera fijado de antemano, sin tener en cuenta que se eligió por
ser la de menor p-valor, y la suma~\eqref{eq:sigtree} supone independencia entre
pruebas que comparten la misma variable $Z$.

\subsection{Relación con la propuesta}\label{sec:sigtree-propuesta}

La propuesta de esta tesis toma de SigTree su principio central, construir una
partición con pruebas de hipótesis y detenerse cuando ninguna división es
significativa, y lo traslada al refinamiento de los bloques de ST-CC. Esa
traslación exige adaptar cada componente. La Tabla~\ref{tab:sigtree-astcc} resume
la correspondencia, que se desarrolla en el Capítulo~\ref{chap:propuesta}.

\begin{table}[ht]
  \centering
  {\footnotesize
  \begin{tabular}{@{}>{\raggedright\arraybackslash}p{0.2\textwidth}
                  >{\raggedright\arraybackslash}p{0.34\textwidth}
                  >{\raggedright\arraybackslash}p{0.36\textwidth}@{}}
    \toprule
    \textbf{Componente} & \textbf{SigTree} & \textbf{Propuesta} \\
    \midrule
    Objeto que se divide & Conjunto de muestras con atributos categóricos &
      Bloque de una matriz de conteos región~$\times$~tiempo \\
    División candidata & Presencia o ausencia de una categoría &
      Bipartición de las filas o de las columnas del bloque \\
    Qué se prueba & Asociación entre la división y los demás atributos &
      Asociación entre la división y la otra dimensión del bloque \\
    Estadístico & Suma de estadísticos $X^2$ & $G^2$ de la tabla colapsada,
      coherente con la I-divergencia \\
    Distribución de referencia & $\chi^2$ asintótica & Simulación condicionada a
      los totales del bloque, repitiendo la búsqueda \\
    Comparaciones múltiples & Bonferroni uniforme & Reparto jerárquico del nivel
      a lo largo del árbol \\
    Criterio de parada & Significancia y reglas triviales & Significancia, tamaño
      de efecto y reglas triviales \\
    \bottomrule
  \end{tabular}
  }
  \caption{Correspondencia entre SigTree \cite{Hu2025} y el método propuesto.}
  \label{tab:sigtree-astcc}
\end{table}

Las limitaciones de los dos métodos base orientan además el diseño. La elección
heurística del número de grupos en ST-CC (Sección~\ref{sec:stcc-limitaciones})
motiva que la resolución final no dependa solo de $K_R$ y $K_T$, y la inflación
del estadístico de SigTree con muestras grandes motiva complementar la
significancia con un umbral de tamaño de efecto.
